\documentclass[11pt,a4paper]{article}
\usepackage[utf8]{inputenc}
\usepackage[T1]{fontenc}
\usepackage[margin=2.4cm]{geometry}
\usepackage{amsmath,amssymb,amsthm,mathtools,booktabs,longtable,array,xcolor,xspace}
\emergencystretch=2em
\usepackage[colorlinks=true,linkcolor=blue!60!black,citecolor=green!40!black,urlcolor=blue!60!black]{hyperref}

\newcommand{\La}{\textsf{\textbf{(a)}}\xspace}
\newcommand{\Lb}{\textsf{\textbf{(b)}}\xspace}
\newcommand{\Lc}{\textsf{\textbf{(c)}}\xspace}
\newcommand{\Ld}{\textsf{\textbf{(d)}}\xspace}
\newcommand{\dd}{\mathrm{d}}
\newcommand{\ii}{\mathrm{i}}
\newcommand{\ee}{\mathrm{e}}
\newcommand{\hb}{\hbar}
\newcommand{\sgn}{\operatorname{sgn}}
\newcommand{\Res}{\operatorname{Res}}
\newcommand{\skin}{\texttt{phase2/\allowbreak scripts/\allowbreak resonant\_\allowbreak channels\_\allowbreak kinematics.py}}
\newcommand{\scons}{\texttt{phase2/\allowbreak scripts/\allowbreak resonant\_\allowbreak channels\_\allowbreak conservation.py}}
\newcommand{\squant}{\texttt{phase2/\allowbreak scripts/\allowbreak resonant\_\allowbreak channels\_\allowbreak quantum.py}}
\newcolumntype{L}[1]{>{\raggedright\arraybackslash}p{#1}}

\newtheorem{proposition}{Proposition}[section]
\newtheorem{theorem}[proposition]{Theorem}
\theoremstyle{remark}
\newtheorem{remark}[proposition]{Remark}

\title{\textbf{The Quantum Fate of the Good Boussinesq Equation's\\
Resonant and Number-Changing Channels}}
\author{Project \emph{quantum\_boussinesq\_solitons} (Ranked next step \#3) --- PI: Gustav Delius}
\date{4 November 2026}

\begin{document}
\maketitle

\begin{abstract}
We investigate ranked next step~\#3 of the Phase~2 master report
(\texttt{phase2/\allowbreak report.pdf}): the classical kinematics, hierarchy
conservation laws, and semiclassical/quantum implications of the resonant
($A_{12}\to 0$ or $\infty$), number-changing ($2\to 1$ merging and $1\to 2$
splitting Y-junctions), and singular ($A_{12}<0$ and $A_{\mathrm{br}}<1/4$)
soliton channels of the \textbf{good} Boussinesq equation
$u_{tt}=u_{xx}-u_{xxxx}-3(u^2)_{xx}$. Every nontrivial claim is labelled \La\
(proved in the literature), \Lb\ (conjectured or heuristic in the literature),
\Lc\ (derived and verified in this project), or \Ld\ (quantum interpretation,
conjecture, or open physical question).

Our main structural results \Lc\ and quantum conclusions \Ld\ are:
\begin{enumerate}
\item \textbf{KP-triangle unification of resonant kinematics and light-branch stability (\skin, \squant).}
In KP root variables, a soliton $(\alpha,\epsilon)\in(0,\pi/2)\times\{\pm 1\}$
carries an ordered chord $(p_a,p_b)$ of the cubic
$4p^3-p=\frac{\ell}{3\sqrt{3}}$ ($\ell=-\epsilon\cos 3\alpha$) on the ellipse
$p_a^2+p_a p_b+p_b^2=\frac{1}{4}$ \La/\Lc. Every nondegenerate real KP triple
$q_1>q_2>q_3$ ($\ell\in(-1,1)$) spans \textbf{two opposite-moving light chords}
$S_{12}=(q_1,q_2)$ ($\alpha_{12}\in(\pi/6,\pi/3),\,\epsilon_{12}=+1$) and
$S_{23}=(q_2,q_3)$ ($\alpha_{23}\in(0,\pi/6),\,\epsilon_{23}=-1$, or with
signs reversed) and \textbf{one heavy chord} $S_{13}=(q_1,q_3)$
($\alpha_{13}=\alpha_{12}+\alpha_{23}+\frac{\pi}{3}\text{ or }\frac{2\pi}{3}-(\alpha_{12}-\alpha_{23})\in(\pi/3,\pi/2)$) \Lc.
The three genuine codimension-one (measure-zero) resonance curves in
$(\alpha_1,\alpha_2)\in(0,\pi/2)^2$---head-on $R_{\mathrm{ho}}^+$
($\alpha_1+\alpha_2=\pi/3,\,A_{\mathrm{ho}}=\infty$), head-on $R_{\mathrm{ho}}^-$
($|\alpha_1-\alpha_2|=\pi/3,\,A_{\mathrm{ho}}=0$), and overtaking
$R_{\mathrm{ov}}$ ($\alpha_1+\alpha_2=2\pi/3,\,A_{\mathrm{ov}}=0$; we confirm
\texttt{audit\_smx.pdf} that $|\alpha_1-\alpha_2|=\pi/3$ is \emph{not} an
overtaking resonance)---are the three pairwise projections of this single
family of KP triangles \Lc. On the light branch $\alpha\in(0,\pi/3)$
($|P|\in(0,\sqrt{3}/8)$), $E_{\mathrm{light}}(P)$ is strictly concave in $|P|$
with $\frac{\dd^2 E_{\mathrm{light}}}{\dd P^2}=-\frac{3}{2\sin 3\alpha}<0$, so
$E_{\mathrm{light}}(P_1)+E_{\mathrm{light}}(P_2)>E_{\mathrm{light}}(P_1+P_2)$
for all nonzero $(P_1,P_2)$ \Lc: \textbf{no light soliton can split into two
solitons ($1_{\mathrm{light}}\not\to 2$), and no two light solitons can fuse
into a light soliton ($2_{\mathrm{light}}\not\to 1_{\mathrm{light}}$)} \Lc.
Conversely, \textbf{every heavy soliton ($\alpha\in(\pi/3,\pi/2)$) is embedded
in the two-light-soliton head-on continuum} with non-vanishing $1\to 2$
phase-space Jacobian $J_{1\to 2}=|v_{12}-v_{23}|^{-1}=(\sqrt{3}\,k_{13})^{-1}$ \Lc.

\item \textbf{Confirmation of the KP-chord hypothesis for classical hierarchy charges, and its one-loop breakdown at fixed physical momentum (\scons).}
Using the hyperbolic polynomial map $T=\tanh(kx/2)$, we integrate the
Gelfand--Dickey densities $\Res L^{n/3}-\Res L_{\mathrm{vac}}^{n/3}$
\textbf{symbolically} on the one-soliton profile through spin $n=8$ ($H_1,H_2,
H_3=0,H_4,H_5,H_6=0,H_7,H_8$, extending Phase~1's numerical fit for $n\le 5$ to
exact symbolic proofs through $n=8$, including the new exact decompositions
$H_7=\Delta_7-\frac{7}{12}\Delta_5+\frac{7}{324}\Delta_1$ and
$H_8=\Delta_8-\frac{7}{36}\Delta_4+\frac{11}{324}\Delta_2$ with
$\Delta_m=(-p_b)^m-(-p_a)^m$) \Lc. Every physical charge ($M,\Pi,P,E$) and every
hierarchy charge $H_n$ has the telescopic chord form $Q(p_a,p_b)=F_Q(p_b)-F_Q(p_a)$
for an explicit polynomial $F_Q(p)$ \Lc. Consequently, chord concatenation
$(q_1,q_2)+(q_2,q_3)\to(q_1,q_3)$ satisfies
$Q(S_{12})+Q(S_{23})=Q(S_{13})$ \textbf{identically for the entire classical
hierarchy} \Lc\ (verified to $10^{-50}$ by direct spatial quadrature on the
exact 3-exponential Y-junction solution across collision). Unlike in
sine-Gordon, the classical conservation laws do \textbf{not} forbid $2\leftrightarrow 1$
fusion \Lc. At one loop, fixed-profile mode-sum shifts $(\Delta E_1^{\mathrm{ren}},
\Delta P_{\mathrm{mode}})$ remain telescopic \Lc, but at fixed physical
(dressed) momentum $P^{\mathrm{phys}}$ the dispersion shift
$F|_P(S)=\Delta E_1^{\mathrm{ren}}(S)-v(S)\Delta P_{\mathrm{mode}}(S)=-L_1(S)$
breaks the chord form because $v(S)=\sqrt{3}(p_a+p_b)$ multiplies
$\Delta P_{\mathrm{mode}}(S)$, yielding a nonzero mismatch
$\Delta(F|_P)=F|_P(S_{12})+F|_P(S_{23})-F|_P(S_{13})=0.38879244070035460274$
at $\phi=0.23$ \Lc.

\item \textbf{Semiclassical phase, unitarity deficit, and three-body factorisation near resonance (\squant).}
At transverse distance $h\to 0^+$ from any of $R_{\mathrm{ov}},R_{\mathrm{ho}}^+,
R_{\mathrm{ho}}^-$, both $\ell_1-\ell_2=c_\ell h+O(h^2)$ and
$A_{12}=c_A h^{\sigma_{\mathrm{res}}}(1+O(h))$ ($\sigma_{\mathrm{res}}=\pm 1$)
have simple zeros/poles, so the classical shift $\Delta x_j=\mp k_j^{-1}\ln A_{12}$
diverges as $\mp\sigma_{\mathrm{res}}k_j^{-1}\ln h$, while the tree-level phase
$\delta(1,2)=T_{12}|_{\mathrm{res}}+\frac{2}{9}\sigma_{\mathrm{res}}c_\ell\,h\ln h+O(h)$
has a finite limit $T_{12}|_{\mathrm{res}}$ with a logarithmic branch-point
slope $\partial_h\delta\propto\ln h$ \Lc\ (confirming \texttt{audit\_smx.pdf}).
At fixed physical momenta, the $O(\hb^0)$ Taylor counterterm
$-\big(\frac{\Delta P_1}{k_1}-\frac{\Delta P_2}{k_2}\big)\ln A_{12}$ diverges as
$\ln h$ on the classical resonance curve \Lc. Continuing into the singular domain
$A_{12}<0$ gives $\operatorname{Im}\delta=\frac{2\pi}{9}|\ell_1-\ell_2|>0$ on
the physical sheet, so $|S_{2\to 2}^{\mathrm{WKB}}|^2=\exp(-\frac{4\pi}{9\hb}|\ell_1-\ell_2|)<1$
violates elastic two-body unitarity \Lc. In 3-soliton collisions with a resonant
pair $S_{12}+S_{23}\to S_{13}$, classical cross-ratios and phases factorise
exactly ($A_{3,12}A_{3,23}=A_{3,13}$ and
$\delta_{\mathrm{cl}}(3,S_{12})+\delta_{\mathrm{cl}}(3,S_{23})=\delta_{\mathrm{cl}}(3,S_{13})$)
and bare one-loop phonon integrals telescope ($I_{SS}(3,S_{12})+I_{SS}(3,S_{23})=I_{SS}(3,S_{13})$),
but the continued Gamma candidate fails the bulk fusion bootstrap by
$\ee^{\mp\ii\frac{\pi}{2}\sgn\eta}(1+O(\hb))$ and the target-soliton momentum
renormalisation fails to telescope \Lc.

\item \textbf{Conclusion on elastic factorised scattering (\Ld\ based on \Lc).}
Purely elastic, unitary, factorised two-soliton scattering is kinematically
protected \textbf{only in the light-light overtaking sector}
($\alpha_1,\alpha_2\in(0,\pi/3),\,\epsilon_1=\epsilon_2$) and in the
sub-threshold **light-light attractive head-on triangle**
($\alpha_1+\alpha_2<\pi/3,\,\epsilon_1=-\epsilon_2$), where $A_{12}>0$ and no
real on-shell $2\to 1$ or $1\to 2$ channel exists \Lc/\Ld. Heavy solitons
($\alpha\in(\pi/3,\pi/2)$) are quantum mechanically unstable resonance states
embedded in the two-light-soliton head-on continuum \Ld, and at or beyond the
resonant boundaries ($A_{12}=0,\infty$ and $A_{12}<0$) elastic 2-body unitarity
must be augmented by the $2\leftrightarrow 1$ fusion/splitting channel \Ld.
\end{enumerate}
\end{abstract}

\tableofcontents
\clearpage

% ==============================================================================
\section{Scope, conventions, and audit precedence}
\label{sec:setup}
% ==============================================================================

We work in the fixed project conventions (\texttt{AGENTS.md},
\texttt{README.md}, \texttt{phase2/\allowbreak report.pdf}) for the
\textbf{good} Boussinesq equation \La:
\begin{equation}\label{eq:good_pde}
  u_{tt} = u_{xx} - u_{xxxx} - 3(u^2)_{xx}, \qquad
  u = \varphi_x = 2(\ln\tau)_{xx}, \qquad
  (D_t^2 - D_x^2 + D_x^4)\,\tau\cdot\tau = 0,
\end{equation}
with Lagrangian density $\mathcal{L} = \frac{1}{2}\varphi_t^2 - \frac{1}{2}\varphi_x^2 + \varphi_x^3 - \frac{1}{2}\varphi_{xx}^2$
and canonical symplectic form $\Omega = \int \dd\pi\wedge\dd\varphi$ ($\pi=\varphi_t$).
A one-soliton with amplitude parameter $\alpha\in(0,\pi/2)$ and direction
$\epsilon\in\{+1,-1\}$ has \La/\Lc:
\begin{align}\label{eq:onesol}
  k &= \sin\alpha, \quad v = \epsilon\cos\alpha, \quad
  \omega = kv = \frac{\epsilon}{2}\sin 2\alpha, \quad
  \ell = -\epsilon\cos 3\alpha, \nonumber\\
  M &= 2k, \quad \Pi = -Mv = -2\epsilon k\cos\alpha, \quad
  P = \frac{2}{3}\epsilon k^3\cos\alpha, \quad
  E = \frac{2}{15}k^3(5-4k^2),
\end{align}
satisfying $\frac{\dd E}{\dd P} = v$ and $\frac{\dd P}{\dd\alpha} = \frac{2}{3}\epsilon k\sin 3\alpha = \frac{2}{9}k\,\ell'(\alpha)$ \Lc.
At the momentum cusp $\alpha = \pi/3$ ($k = \sqrt{3}/2$, $|v| = 1/2$, $|\ell|=1$),
$|P|$ attains its global maximum $P_{\max} = \sqrt{3}/8 \approx 0.21650635$ and
$\frac{\dd P}{\dd\alpha} = 0$, dividing the one-soliton manifold into the
\textbf{light branch} $\alpha\in(0,\pi/3)$ ($k<\sqrt{3}/2$, $|v|>1/2$) and the
\textbf{heavy branch} $\alpha\in(\pi/3,\pi/2)$ ($k>\sqrt{3}/2$, $|v|<1/2$) \Lc\
(\texttt{phase1/\allowbreak review.pdf}).

\paragraph{Audit precedence.}
Following the ground rules of the task, audit verdicts override package claims.
In particular, we incorporate the corrections from \texttt{audit\_smx.pdf},
\texttt{audit\_defq.pdf}, \texttt{audit\_det\_projected.pdf}, and
\texttt{errata.pdf}:
\begin{itemize}
\item The line $|\alpha_1-\alpha_2|=\pi/3$ is \emph{not} an overtaking resonance
(\texttt{audit\_smx.pdf}); we verify this directly in \skin.
\item The tree-level S-matrix limit at a genuine resonance is continuous, not
differentiably ($C^1$) smooth (\texttt{audit\_smx.pdf}); we derive the exact
$h\ln h$ coefficient on all three genuine resonance lines in \squant.
\item Every two-phase zero-background Hirota breather of the good equation has
$A_{\mathrm{br}} = \frac{1-3\cot^2 a}{1+3\coth^2 b} < \frac{1}{4} < 1$ and is
singular (\texttt{phase1/\allowbreak review.pdf}, \texttt{audit\_smx.pdf}); this
does not by itself exclude quantum bound states beyond the two-phase Hirota
ansatz \Ld.
\item In \texttt{audit\_defq.pdf}, the classical KP-edge charge factorisation
survives, while the canonical-charge triangle holonomy no-go theorem is
qualified as conditional on fixed unshifted absorption rapidities and single-valued
vertex counterterms.
\end{itemize}

% ==============================================================================
\section{Question 1: Classical kinematics of resonant and number-changing channels}
\label{sec:kinematics}
% ==============================================================================

\subsection{KP root/chord dictionary and the three genuine resonance loci}

In the third-order Lax representation $L = \partial_x^3 + (\frac{3}{2}u - \frac{1}{4})\partial_x + r$
\La~\cite{Zak,Dickey}, a soliton $(\alpha,\epsilon)$ corresponds to an ordered
pair (chord) $(p_a,p_b)$ ($p_a > p_b$) of real roots of the KP spectral cubic \La/\Lc\
(\texttt{phase1/\allowbreak review.pdf}, \texttt{errata.pdf} M06--M07, \skin):
\begin{equation}\label{eq:kp_roots}
  p_a = \frac{1}{2}\!\left(\frac{v}{\sqrt{3}} + k\right)
      = \frac{\epsilon}{\sqrt{3}}\cos\!\left(\alpha - \frac{\epsilon\pi}{3}\right), \qquad
  p_b = \frac{1}{2}\!\left(\frac{v}{\sqrt{3}} - k\right)
      = \frac{\epsilon}{\sqrt{3}}\cos\!\left(\alpha + \frac{\epsilon\pi}{3}\right),
\end{equation}
so that $p_a - p_b = k > 0$, $p_a + p_b = \frac{v}{\sqrt{3}}$, and both roots
satisfy
\begin{equation}\label{eq:kp_cubic}
  4p^3 - p = \frac{\ell}{3\sqrt{3}} = -\frac{\epsilon}{3\sqrt{3}}\cos 3\alpha, \qquad
  p_a^2 + p_a p_b + p_b^2 = \frac{1}{4}.
\end{equation}
The Hirota two-soliton interaction coefficient
$A_{12} = -\frac{\mathcal{P}(k_1-k_2,\omega_1-\omega_2)}{\mathcal{P}(k_1+k_2,\omega_1+\omega_2)}$
($\mathcal{P}(K,W) = W^2 - K^2 + K^4$) equals the KP cross-ratio in both
relative orientations \La~\cite{DMHC}, \Lc\ (\skin):
\begin{equation}\label{eq:A12_crossratio}
  A_{12} = \frac{(p_{1a} - p_{2a})(p_{1b} - p_{2b})}{(p_{1a} - p_{2b})(p_{1b} - p_{2a})}
  = \begin{cases}
    \dfrac{1 - 3\cot^2\!\big(\frac{\alpha_1+\alpha_2}{2}\big)}{1 - 3\cot^2\!\big(\frac{\alpha_1-\alpha_2}{2}\big)}
    & (\text{overtaking, }\epsilon_1 = \epsilon_2), \\[10pt]
    \dfrac{1 - 3\tan^2\!\big(\frac{\alpha_1-\alpha_2}{2}\big)}{1 - 3\tan^2\!\big(\frac{\alpha_1+\alpha_2}{2}\big)}
    & (\text{head-on, }\epsilon_1 = -\epsilon_2).
  \end{cases}
\end{equation}

\begin{proposition}[Exact resonance loci and KP-triangle unification, \Lc, \skin]\label{prop:resonances}
Let $(\alpha_1,\alpha_2) \in Q := (0,\pi/2)^2$ with $\alpha_1 \ne \alpha_2$ when
$\epsilon_1 = \epsilon_2$.
\begin{enumerate}
\item \textbf{Characterization by shared KP roots and equal $\ell$:}
A resonant interaction $A_{12} = 0$ or $A_{12} = \infty$ occurs if and only if
solitons $1$ and $2$ share a root of the KP cubic \eqref{eq:kp_cubic}, which in
$Q = (0,\pi/2)^2$ implies $\ell_1 = \ell_2$ and defines three real-analytic
codimension-one lines:
\begin{itemize}
\item \emph{Overtaking resonance $R_{\mathrm{ov}}$ ($\epsilon_1 = \epsilon_2$, $A_{\mathrm{ov}} = 0$):}
\begin{equation}\label{eq:R_ov}
  \alpha_1 + \alpha_2 = \frac{2\pi}{3}, \qquad \alpha_1,\alpha_2 \in \left(\frac{\pi}{6},\frac{\pi}{2}\right),
\end{equation}
where $p_{1a} = p_{2a}$ (for $\epsilon_1=\epsilon_2=+1$) or $p_{1b} = p_{2b}$
(for $\epsilon_1=\epsilon_2=-1$), and $\ell_1 = \ell_2 = -\epsilon_1\cos 3\alpha_1 > 0$.
The overtaking denominator $1 - 3\cot^2\!\big(\frac{\alpha_1-\alpha_2}{2}\big)$
vanishes only at $|\alpha_1-\alpha_2| = \frac{2\pi}{3} \notin Q$, so
$A_{\mathrm{ov}}$ has \textbf{no pole} in $Q$. In particular, on the false line
$|\alpha_1-\alpha_2| = \frac{\pi}{3}$ (refuted in \texttt{audit\_smx.pdf}), the
denominator equals $1 - 3\cot^2(\pi/6) = -8 \ne 0$; at $(\alpha_1,\alpha_2)=(0.2,0.2+\pi/3)$
we have $A_{\mathrm{ov}} = 0.35546547984906164962$ and
$\ell_1-\ell_2 = -2\cos(0.6) = -1.6506712298193565945 \ne 0$.

\item \emph{Head-on attractive resonance $R_{\mathrm{ho}}^+$ ($\epsilon_1 = -\epsilon_2$, $A_{\mathrm{ho}} = \infty$):}
\begin{equation}\label{eq:R_ho_plus}
  \alpha_1 + \alpha_2 = \frac{\pi}{3}, \qquad \alpha_1,\alpha_2 \in \left(0,\frac{\pi}{3}\right),
\end{equation}
where $p_{1b} = p_{2a}$ (for $\epsilon_1=+1,\epsilon_2=-1$) or $p_{1a} = p_{2b}$
(for $\epsilon_1=-1,\epsilon_2=+1$), and $\ell_1 = \ell_2 = -\epsilon_1\cos 3\alpha_1$.

\item \emph{Head-on repulsive resonance $R_{\mathrm{ho}}^-$ ($\epsilon_1 = -\epsilon_2$, $A_{\mathrm{ho}} = 0$):}
\begin{equation}\label{eq:R_ho_minus}
  |\alpha_1 - \alpha_2| = \frac{\pi}{3}, \qquad
  \min(\alpha_1,\alpha_2) \in \left(0,\frac{\pi}{6}\right), \quad
  \max(\alpha_1,\alpha_2) \in \left(\frac{\pi}{3},\frac{\pi}{2}\right),
\end{equation}
where $p_{1b} = p_{2b}$ (for $\alpha_1 - \alpha_2 = \pi/3,\,\epsilon_1=+1,\epsilon_2=-1$)
or $p_{1a} = p_{2a}$ (for $\alpha_2 - \alpha_1 = \pi/3,\,\epsilon_1=+1,\epsilon_2=-1$),
and $\ell_1 = \ell_2$ because $\cos(3\alpha_2\pm\pi) = -\cos 3\alpha_2$
compensates $\epsilon_1 = -\epsilon_2$.
\end{itemize}

\item \textbf{Unification via KP triangles and light/heavy branch assignment:}
Every nondegenerate KP cubic $4q^3 - q = \frac{\ell}{3\sqrt{3}}$ with
$\ell \in (-1,1)$ has three real roots ordered $q_1 > q_2 > q_3$, parametrised
by $\phi \in (0,\pi/3)$ as
\begin{align}\label{eq:kp_triple_param}
  q_1 &= \frac{1}{\sqrt{3}}\cos\phi \in \left(\frac{1}{2\sqrt{3}},\frac{1}{\sqrt{3}}\right), \qquad
  q_2 = \frac{1}{\sqrt{3}}\cos\!\left(\phi-\frac{2\pi}{3}\right) \in \left(-\frac{1}{2\sqrt{3}},\frac{1}{2\sqrt{3}}\right), \nonumber\\
  q_3 &= \frac{1}{\sqrt{3}}\cos\!\left(\phi+\frac{2\pi}{3}\right) \in \left(-\frac{1}{\sqrt{3}},-\frac{1}{2\sqrt{3}}\right).
\end{align}
The three chords $S_{12}=(q_1,q_2)$, $S_{23}=(q_2,q_3)$, and $S_{13}=(q_1,q_3)$
have wavenumbers and velocities:
\begin{align}
  S_{12}&:\ k_{12} = q_1-q_2 = \sin\!\left(\frac{\pi}{3}-\phi\right) < \frac{\sqrt{3}}{2}, &
  v_{12} &= -\sqrt{3}\,q_3 = +\cos\!\left(\frac{\pi}{3}-\phi\right) > \frac{1}{2}
  \ \ (\text{right \textbf{light}}), \nonumber\\
  S_{23}&:\ k_{23} = q_2-q_3 = \sin\phi < \frac{\sqrt{3}}{2}, &
  v_{23} &= -\sqrt{3}\,q_1 = -\cos\phi < -\frac{1}{2}
  \ \ (\text{left \textbf{light}}), \label{eq:triple_chords}\\
  S_{13}&:\ k_{13} = q_1-q_3 = \sin\!\left(\phi+\frac{\pi}{3}\right) > \frac{\sqrt{3}}{2}, &
  v_{13} &= -\sqrt{3}\,q_2 = \cos\!\left(\phi+\frac{\pi}{3}\right) \in \left(-\frac{1}{2},\frac{1}{2}\right)
  \ \ (\textbf{heavy},\ \phi\ne\tfrac{\pi}{6}). \nonumber
\end{align}
For $\phi \in (0,\pi/6)$ ($\ell > 0$), $S_{13}$ is a right-moving heavy soliton
with $\alpha_{13} = \phi+\frac{\pi}{3} \in (\pi/3,\pi/2)$, and the three pairs
of chords of the same KP triangle $\{S_{12},S_{23},S_{13}\}$ project precisely
onto the three resonance lines:
\begin{equation}\label{eq:three_projections}
  (S_{12},S_{23}) \in R_{\mathrm{ho}}^+ \ (\text{light--light}), \quad
  (S_{13},S_{23}) \in R_{\mathrm{ho}}^- \ (\text{heavy--light}), \quad
  (S_{12},S_{13}) \in R_{\mathrm{ov}} \ (\text{light--heavy}).
\end{equation}
(For $\phi \in (\pi/6,\pi/3)$, $\ell < 0$ and $S_{13}$ is a left-moving heavy
soliton with $\alpha_{13} = \frac{2\pi}{3}-\phi \in (\pi/3,\pi/2)$, giving the
mirror-image projections.)
\end{enumerate}
\end{proposition}

\subsection{Reachability from real asymptotic states and domain measure}

On any of the three resonance curves, dropping the vanishing or divergent
four-exponential term (after shifting $\theta_j$ when $A_{12}\to\infty$) yields
the exact three-exponential Hirota Y-junction \La~\cite{HI,LMK,RS,DMHC,CZ}:
\begin{equation}\label{eq:tau_Y}
  \tau_Y(x,t) = c_1\,\ee^{\eta_1} + c_2\,\ee^{\eta_2} + c_3\,\ee^{\eta_3}, \qquad
  \eta_i = \pm\big(q_i x - \sqrt{3}\,q_i^2 t\big) + \eta_{i,0} \quad (c_i > 0).
\end{equation}
Because $\mathcal{P}\big(q_i-q_j,\sqrt{3}(q_i^2-q_j^2)\big) = 4(q_i-q_j)^2\big(q_i^2+q_i q_j+q_j^2-\frac{1}{4}\big) = 0$
for every pair $(i,j)$ of a KP triple \eqref{eq:kp_triple_param}, $\tau_Y(x,t)>0$
solves $(D_t^2-D_x^2+D_x^4)\tau_Y\cdot\tau_Y = 0$ identically and $u_Y = 2(\ln\tau_Y)_{xx}$
is globally smooth for all $(x,t)\in\mathbb{R}^2$ \La/\Lc\ (verified to $10^{-50}$
against the nonlinear PDE in \skin).

\begin{proposition}[Asymptotic reachability and measure of resonant/singular domains, \Lc, \skin]\label{prop:reachability}
Let $q_1 > q_2 > q_3$ be a nondegenerate KP triple with velocities
$v_{12} > v_{13} > v_{23}$ as in \eqref{eq:triple_chords}.
\begin{enumerate}
\item \textbf{Exact $1\to 2$ splitting and $2\to 1$ merging asymptotics:}
Along the ray $x = v_{13} t + z$, $\eta_1 - \eta_3 = \pm(q_1-q_3)z$ is
stationary while
\begin{equation}\label{eq:transverse_exponent}
  \eta_2 - \eta_1 = \pm\sqrt{3}\,(q_1 - q_2)(q_2 - q_3)\,t + O(z)
                  = \pm\sqrt{3}\,k_{12}k_{23}\,t + O(z).
\end{equation}
Since $k_{12}k_{23} > 0$:
\begin{itemize}
\item For the $+$ sign in \eqref{eq:tau_Y}, as $t\to -\infty$ the middle
exponential $\ee^{\eta_2}$ is suppressed by $\ee^{-\sqrt{3}k_{12}k_{23}|t|}$ and
$u_Y(x,t)$ approaches a \textbf{single heavy soliton} $S_{13}$ moving along
$x = v_{13}t$; as $t\to +\infty$, $u_Y(x,t)$ splits ($1\to 2$) into the
\textbf{two well-separated light solitons} $S_{12}$ (moving right at $v_{12}>1/2$)
and $S_{23}$ (moving left at $v_{23}<-1/2$).
\item For the $-$ sign in \eqref{eq:tau_Y} (time reversal $t\mapsto -t,x\mapsto -x$),
as $t\to -\infty$ the solution consists of \textbf{two real, smooth,
exponentially separated incoming light solitons} $S_{12}$ (at $x\sim v_{12}t<0$)
and $S_{23}$ (at $x\sim v_{23}t>0$) colliding head-on and \textbf{merging
($2\to 1$)} into the single heavy soliton $S_{13}$ as $t\to +\infty$ (verified
at $t=\pm 400$ to $10^{-48}$ in \skin).
\end{itemize}
Thus the resonant $2\to 1$ and $1\to 2$ channels are \textbf{directly reachable
from real, nonsingular asymptotic soliton states}. What appears as an overtaking
resonance $R_{\mathrm{ov}}$ ($S_{12}+S_{13}\leftrightarrow S_{23}$) or repulsive
head-on resonance $R_{\mathrm{ho}}^-$ ($S_{13}+S_{23}\leftrightarrow S_{12}$) in
two-soliton coordinates is the \emph{same} spatial Y-junction viewed by pairing
the heavy leg $S_{13}$ with one of the two light legs.

\item \textbf{Measure zero of resonant loci vs.\ positive measure of singular regions:}
\begin{itemize}
\item The resonant loci $R_{\mathrm{ov}}$, $R_{\mathrm{ho}}^+$, and
$R_{\mathrm{ho}}^-$ are smooth 1D curves in the 2D kinematic square
$Q = (0,\pi/2)^2$ (and in the four $(\text{light/heavy})\times(\text{light/heavy})$
sheets over $(|P_1|,|P_2|)\in(0,\sqrt{3}/8)^2$), hence of \textbf{two-body
Lebesgue measure zero} (codimension $1$).
\item However, in the \textbf{one-body} kinematic domain of the heavy branch
$\alpha_{13}\in(\pi/3,\pi/2)$, the $1\to 2$ decay channel $S_{13}\to S_{12}+S_{23}$
is open on the \textbf{entire heavy branch} (full 1D measure, codimension $0$)!
\item In the 2D two-body square $Q=(0,\pi/2)^2$ (area $\pi^2/4$), the open
singular regions ($A_{12}<0$, where the 4-exponential Hirota solution has
logarithmic poles $\tau=0$) have \textbf{positive area}: $\frac{2}{9}$ of $Q$
for overtaking ($\alpha_1+\alpha_2>2\pi/3$) and $\frac{2}{3}$ of $Q$ for head-on
($\alpha_1+\alpha_2>\pi/3$ with $|\alpha_1-\alpha_2|<\pi/3$). Crucially, inside
the \textbf{light-light square} $(0,\pi/3)^2$ (area $\pi^2/9$, $\frac{4}{9}$ of
$Q$), \textbf{$100\%$ of overtaking collisions} ($\alpha_1+\alpha_2<2\pi/3$)
and \textbf{$50\%$ of head-on collisions} ($\alpha_1+\alpha_2<\pi/3$) are
strictly regular ($A_{12}>0$) and non-resonant \Lc.
\end{itemize}
\end{enumerate}
\end{proposition}

% ==============================================================================
\section{Question 2: Conservation laws and the KP-chord hypothesis}
\label{sec:conservation}
% ==============================================================================

\subsection{Exact symbolic proof of the chord form through spin \texorpdfstring{$8$}{8}}

In the Gelfand--Dickey hierarchy \La~\cite{Dickey}, the conserved charges are
$H_n = \int_{-\infty}^\infty \big(\Res L^{n/3} - \Res L_{\mathrm{vac}}^{n/3}\big)\,\dd x$,
where $L = \partial_x^3 + q\,\partial_x + r$ has $q = \frac{3}{2}u - \frac{1}{4}$
and $r = \frac{1}{2}q_x - \frac{v}{2\sqrt{3}}(q+\frac{1}{4})$ on a soliton
moving at velocity $v$. For $n = 3m$, $L^m$ is a purely differential operator so
$\Res L^m \equiv 0$ identically \La. In Phase~1 (\texttt{phase1/\allowbreak scripts/\allowbreak charges.py}),
$H_n$ was integrated numerically for $n \in \{1,2,4,5\}$ and fitted to linear
combinations of $\Delta_m := \mu_1^m - \mu_2^m = (-p_b)^m - (-p_a)^m$. Here we
test the KP-chord hypothesis \textbf{symbolically} through $n = 8$ by mapping
the one-soliton profile to a polynomial in $T = \tanh(kx/2) \in (-1,1)$ \Lc\ (\scons).

\begin{theorem}[Symbolic chord-potential structure of the Boussinesq hierarchy, \Lc, \scons]\label{thm:chord_charges}
Let $(p_a,p_b)$ ($p_a > p_b$) be the KP chord of a soliton $(\alpha,\epsilon)$ on
$p_a^2 + p_a p_b + p_b^2 = \frac{1}{4}$, and define
$\Delta_m := (-p_b)^m - (-p_a)^m = (-1)^{m-1}(p_a^m - p_b^m)$.
\begin{enumerate}
\item \textbf{Exact symbolic integration of $H_1,\dots,H_8$:}
Under $T = \tanh(kx/2)$ ($u = \frac{k^2}{2}(1-T^2)$,
$\partial_x = \frac{k}{2}(1-T^2)\frac{\dd}{\dd T}$,
$\dd x = \frac{2}{k(1-T^2)}\,\dd T$), the vacuum-subtracted densities
$\Res L^{n/3} - \Res L_{\mathrm{vac}}^{n/3}$ reduce to exact polynomials in $T$
whose integrals over $T \in [-1,1]$ decompose modulo
$p_a^2 + p_a p_b + p_b^2 - \frac{1}{4} = 0$ as:
\begin{align}
  H_1 &= \Delta_1, &
  H_2 &= \Delta_2, &
  H_3 &= 0, \nonumber\\
  H_4 &= \Delta_4 - \frac{1}{3}\Delta_2, &
  H_5 &= \Delta_5 - \frac{5}{72}\Delta_1, &
  H_6 &= 0, \label{eq:H1_to_H8}\\
  H_7 &= \Delta_7 - \frac{7}{12}\Delta_5 + \frac{7}{324}\Delta_1, &
  H_8 &= \Delta_8 - \frac{7}{36}\Delta_4 + \frac{11}{324}\Delta_2. \nonumber
\end{align}
\item \textbf{Univariate chord potentials $Q(p_a,p_b) = F_Q(p_b) - F_Q(p_a)$:}
Every hierarchy charge $H_n$ ($1\le n\le 8$) and every physical charge
$M,\Pi,P,E$ has the exact chord-difference form $Q(p_a,p_b) = F_Q(p_b) - F_Q(p_a)$
modulo $p_a^2 + p_a p_b + p_b^2 = \frac{1}{4}$, with univariate polynomials:
\begin{align}
  F_M(p) &= -2p, &
  F_\Pi(p) &= 2\sqrt{3}\,p^2, \nonumber\\
  F_P(p) &= -2\sqrt{3}\,p^4 + \frac{1}{\sqrt{3}}\,p^2, &
  F_E(p) &= -\frac{24}{5}\,p^5 + \frac{1}{6}\,p, \label{eq:F_potentials}\\
  F_{H_7}(p) &= -p^7 + \frac{7}{12}\,p^5 - \frac{7}{324}\,p, &
  F_{H_8}(p) &= p^8 - \frac{7}{36}\,p^4 + \frac{11}{324}\,p^2. \nonumber
\end{align}
\item \textbf{Why all hierarchy charges have the chord form, and the spin-$3m$ gap:}
On the zero-background KP hierarchy ($\tilde{L} = \partial_x^3 + \tilde{q}\partial_x + \tilde{r}$,
$\tilde{q}_\infty = 0$), a line soliton with chord $(\mu_1,\mu_2)$ is generated
by a Sato dressing transformation shifting $\ln\tau$ by $\ln\frac{\partial-\mu_1}{\partial-\mu_2}$,
so $\int\Res\tilde{L}^{n/3}\,\dd x = \mu_1^n - \mu_2^n = \Delta_n$ \La~\cite{Dickey}.
Shifting to the good Boussinesq constant background $L = \tilde{L} - \frac{1}{4}\partial_x$
mixes $\tilde{L}^{n/3}$ only with lower powers $\tilde{L}^{(n-2j)/3}$ of the
same parity, so every hierarchy charge $H_n$ on a one-soliton is a linear
combination of $\Delta_m$ ($m\le n,\,m\equiv n\pmod 2$) and therefore has the
chord form $F_{H_n}(p_b) - F_{H_n}(p_a)$ \La/\Lc.
Crucially, however, \emph{not every} polynomial in $(k,v)$ is a chord difference
of matching degree: because $\Delta_3 = \frac{1}{4}\Delta_1$ on
$p_a^2+p_a p_b+p_b^2=\frac{1}{4}$, the space $\operatorname{span}\{\Delta_1,\Delta_5\}$
contains $k^3(5-4k^2) = 36\big(\Delta_5 - \frac{5}{144}\Delta_1\big)$ (the exact
functional form of $E$!), but does \textbf{not} contain pure $k^3$ \Lc\ (\scons).
\end{enumerate}
\end{theorem}

\subsection{Telescoping in \texorpdfstring{$2\leftrightarrow 1$}{2<->1} fusion and one-loop fate at fixed momentum}

\begin{proposition}[Classical conservation vs.\ one-loop fixed-momentum non-telescoping, \Lc, \scons]\label{prop:fusion_conservation}
Let $q_1 > q_2 > q_3$ be any KP triple with chords $S_{12}=(q_1,q_2)$,
$S_{23}=(q_2,q_3)$, and $S_{13}=(q_1,q_3)$.
\begin{enumerate}
\item \textbf{Exact classical conservation of the infinite hierarchy:}
For every charge $Q$ of the chord form $Q(p_a,p_b) = F_Q(p_b) - F_Q(p_a)$
(including $M,\Pi,P,E$ and all $H_n$), chord concatenation $(q_1,q_2)+(q_2,q_3)\to(q_1,q_3)$
telescopes identically:
\begin{equation}\label{eq:chord_telescope}
  Q(S_{12}) + Q(S_{23}) = \big(F_Q(q_2) - F_Q(q_1)\big) + \big(F_Q(q_3) - F_Q(q_2)\big)
  = F_Q(q_3) - F_Q(q_1) = Q(S_{13}).
\end{equation}
Direct 60-digit spatial quadrature of the nonlinear charge densities of $M,\Pi,P,E$
on the exact 3-exponential Y-junction solution $u_Y(x,t)$ at $t \in \{-2,0,+2\}$
agrees with $Q(S_{12})+Q(S_{23}) = Q(S_{13})$ to $< 10^{-50}$ \Lc.
Thus \textbf{the classical conservation laws do not forbid $2\leftrightarrow 1$
soliton fusion or splitting}, in sharp contrast to sine-Gordon (where kinks carry
a fixed topological charge and a 1-parameter dispersion $Q_s = m^s\sinh(s\theta)$,
so $Q_s(\theta_1)+Q_s(\theta_2)=Q_s(\theta_3)$ has no real solution).

\item \textbf{Equivalence to the defect KP-triangle holonomy and one-loop breakdown at fixed $P^{\mathrm{phys}}$:}
For a Corrigan--Zambon defect of orientation $\eta\in\{\pm 1\}$ on the KP
triangle $\{q_1,q_2,q_3\}$ ($q_1>q_2>q_3$), the absorption signs on the three
directed edges $(q_1\to q_2,q_2\to q_3,q_3\to q_1)$ are
$\rho = (+\eta,+\eta,-\eta)$ (\texttt{defq\_charges\_fusion.pdf},
\texttt{audit\_defq.pdf}). Hence the directed KP-triangle holonomy of any
soliton observable $\mathcal{O}$ is literally the $2\to 1$ fusion balance:
\begin{equation}\label{eq:holonomy_is_fusion}
  \sum_{i=1}^3 \rho_i\,\mathcal{O}(S_i)
  = \eta\,\Big(\mathcal{O}(S_{12}) + \mathcal{O}(S_{23}) - \mathcal{O}(S_{13})\Big).
\end{equation}
Consequently:
\begin{itemize}
\item In the \emph{fixed-profile mode-sum} one-loop prescription,
$\Delta E_1^{\mathrm{ren}}(S)$ and $\Delta P_{\mathrm{mode}}(S)$ are 2-root
dipole differences on the KP cubic and telescope to zero
($\Delta E_1^{\mathrm{ren}}(S_{12})+\Delta E_1^{\mathrm{ren}}(S_{23})-\Delta E_1^{\mathrm{ren}}(S_{13}) = 0$
to $10^{-61}$ at $\phi=0.23$) \Lc.
\item However, in the \emph{fixed-physical-momentum} expansion
$E^{\mathrm{phys}}(P^{\mathrm{phys}}) = E_{\mathrm{cl}}(P^{\mathrm{phys}}) + \hb\,F|_P(P^{\mathrm{phys}}) + O(\hb^2)$,
the charge-pair-invariant one-loop shift
$F|_P(S) = \Delta E_1^{\mathrm{ren}}(S) - v(S)\,\Delta P_{\mathrm{mode}}(S) = -L_1(S)$
(\texttt{audit\_det\_projected.pdf}) is \textbf{not} a chord difference because
$v(S) = \sqrt{3}(p_a+p_b)$ multiplies $\Delta P_{\mathrm{mode}}(S)$:
\begin{align}\label{eq:delta_F_at_P}
  \Delta(F|_P) &:= F|_P(S_{12}) + F|_P(S_{23}) - F|_P(S_{13}) \nonumber\\
  &= -(v_{12}-v_{13})\,\Delta P_{\mathrm{mode}}(S_{12}) - (v_{23}-v_{13})\,\Delta P_{\mathrm{mode}}(S_{23}) \ne 0.
\end{align}
At $\phi = 0.23$, $\Delta(F|_P) = +0.38879244070035460274 \ne 0$ \Lc\ (\scons).
\end{itemize}
\end{enumerate}
\end{proposition}

\begin{remark}[Interpretation for quantum fusion, \Ld]\label{rem:quantum_conservation}
For a single charge pair $(E^{\mathrm{phys}},P^{\mathrm{phys}})$, a nonzero
one-loop shift $\Delta(F|_P) \ne 0$ does \emph{not} by itself forbid $2\to 1$
fusion: because $v_{12} \ne v_{13}$ and $v_{23} \ne v_{13}$, shifting the
resonant momenta by $P_j^{\mathrm{phys}} = P_j^{\mathrm{cl}} + \hb\,\delta P_j$
with $(v_{12}-v_{13})\delta P_{12} + (v_{23}-v_{13})\delta P_{23} = -\Delta(F|_P)$
(for instance $\delta P_j = \Delta P_{\mathrm{mode}}(S_j)$, corresponding to
unshifted bare KP roots) restores $E^{\mathrm{phys}}(P_1^{\mathrm{phys}}) + E^{\mathrm{phys}}(P_2^{\mathrm{phys}}) = E^{\mathrm{phys}}(P_1^{\mathrm{phys}}+P_2^{\mathrm{phys}})$
to $O(\hb)$ along a shifted 1D curve \Lc/\Ld. However, for the \emph{entire
infinite tower} of quantum charges $\{\hat{H}_n\}$ to allow $2\to 1$ fusion on
the \emph{same} shifted curve in $(P_1^{\mathrm{phys}},P_2^{\mathrm{phys}})$, all
one-loop charge corrections $\Delta H_n|_P$ would have to be absorbed by a
single renormalisation of the KP triple $\{q_1,q_2,q_3\}$. Whether the quantum
Boussinesq hierarchy preserves this exact KP-chord elasticity at one loop and
beyond remains open \Ld.
\end{remark}

% ==============================================================================
\section{Question 3: Quantum and semiclassical implications}
\label{sec:quantum}
% ==============================================================================

\subsection{Semiclassical phase \texorpdfstring{$\delta(1,2)$}{delta(1,2)} and divergent time delay as \texorpdfstring{$A_{12}\to 0$ or $\infty$}{A12 -> 0 or infty}}

Recall the audited Darboux-frame semiclassical two-soliton phase \Lc\
(\texttt{phase2/\allowbreak report.pdf} Prop.~2.3, \texttt{audit\_smx.pdf}):
\begin{equation}\label{eq:delta_semi}
  \delta(1,2) = \frac{2}{9}(\ell_1 - \ell_2)\ln A_{12} + T_{12}, \qquad
  T_{12} = \frac{2}{3}M_1 M_2(v_1 - v_2) = -\frac{2}{3}(\Pi_1 M_2 - \Pi_2 M_1),
\end{equation}
whose canonical momentum derivatives give the Darboux shifts
$\Delta X_1' = -\frac{\partial\delta}{\partial P_1} = -k_1^{-1}\ln A_{12} + \gamma_1 M_2$
($\gamma_j = \frac{\dd\Pi_j/\dd\alpha_j}{\dd P_j/\dd\alpha_j}$).

\begin{proposition}[Transverse $h\ln h$ branch point and divergent delay at resonance, \Lc, \squant]\label{prop:phase_resonance}
Let $h \to 0^+$ be the transverse rapidity displacement into the regular domain
($A_{12} > 0$) from any of the three genuine resonance lines $R_{\mathrm{ov}}$
($\alpha_2 = \frac{2\pi}{3}-\alpha_1-h$, $\sigma_{\mathrm{res}}=+1$),
$R_{\mathrm{ho}}^+$ ($\alpha_2 = \frac{\pi}{3}-\alpha_1-h$, $\sigma_{\mathrm{res}}=-1$),
or $R_{\mathrm{ho}}^-$ ($\alpha_1 = \alpha_2+\frac{\pi}{3}+h$, $\sigma_{\mathrm{res}}=+1$).
\begin{enumerate}
\item \textbf{Simple zero/pole of $A_{12}$ and linear vanishing of $\ell_1-\ell_2$:}
On all three lines, one KP root difference $p_{1r}-p_{2s} = O(h)$ vanishes to
first order, so both $\ell_1-\ell_2$ and $A_{12}^{\sigma_{\mathrm{res}}}$ vanish
linearly in $h$:
\begin{equation}\label{eq:transverse_asympt}
  \ell_1 - \ell_2 = c_\ell(\alpha_*)\,h + O(h^2), \qquad
  A_{12}(h) = c_A(\alpha_*)\,h^{\sigma_{\mathrm{res}}}\big(1 + O(h)\big) \qquad (c_A > 0),
\end{equation}
with $c_\ell(\alpha_*) = -3\sin 3\alpha_* \ne 0$ for $\alpha_* \ne \pi/3$.
Consequently, as $h \to 0^+$:
\begin{equation}\label{eq:delta_hlnh}
  \boxed{\;
  \delta(1,2) = T_{12}\big|_{\mathrm{res}}
  + \frac{2}{9}\sigma_{\mathrm{res}}\,c_\ell(\alpha_*)\,h\ln h + O(h), \qquad
  \Delta x_1 = -\frac{\sigma_{\mathrm{res}}}{k_1}\ln h - \frac{\ln c_A}{k_1} + O(h).
  \;}
\end{equation}
Thus $\delta(1,2)$ approaches the finite topological limit $T_{12}|_{\mathrm{res}}$
continuously, while its transverse slope $\partial_h\delta \sim \frac{2}{9}\sigma_{\mathrm{res}}c_\ell\ln h$
and the classical position shift $\Delta x_j$ (and Wigner time delay
$\Delta t_{12} = \frac{\Delta x_1 - \Delta x_2}{v_1 - v_2}$) \textbf{diverge
logarithmically as $\ln(1/h)$} \Lc.

\item \textbf{Logarithmic divergence of the fixed-physical-momentum $O(\hb^0)$ phase:}
While the bare fixed-rapidity one-loop phonon integral $I_{12}(\alpha_1,\alpha_2)$
remains finite on all three resonance lines (\texttt{oneloop\_phase.pdf},
\texttt{audit\_smx.pdf}), expressing the S-matrix at fixed physical momenta
$P_j^{\mathrm{phys}} = P_j^{\mathrm{cl}} + \hb\,\Delta P_j$ introduces the
$O(\hb^0)$ chain-rule counterterm
\begin{equation}\label{eq:delta_renorm_div}
  \delta_{12,\mathrm{renorm}}
  = -\frac{\partial\delta_{\mathrm{cl}}}{\partial P_1}\Delta P_1
    -\frac{\partial\delta_{\mathrm{cl}}}{\partial P_2}\Delta P_2
  = -\left(\frac{\Delta P_1}{k_1} - \frac{\Delta P_2}{k_2}\right)\ln A_{12}
    + \gamma_1 M_2\Delta P_1 - \gamma_2 M_1\Delta P_2,
\end{equation}
whose coefficient $-\big(\frac{\Delta P_1}{k_1}-\frac{\Delta P_2}{k_2}\big)$ is
generically non-zero on all three resonance lines, producing a \textbf{logarithmic
divergence $\propto \ln h$ at $O(\hb^0)$} \Lc:
\begin{center}\footnotesize
\begin{tabular}{@{}lrrrr@{}}\toprule
Resonance point & $T_{12}|_{\mathrm{res}}$ & $c_A$ & $\frac{2}{9}\sigma_{\mathrm{res}}c_\ell$ ($h\ln h$) & $-\big(\frac{\Delta P_1}{k_1}-\frac{\Delta P_2}{k_2}\big)$ ($\ln A_{12}$) \\\midrule
$R_{\mathrm{ov}}$ ($\alpha_1=0.75$) & $+0.898441539345$ & $0.074538040233$ & $-0.518715464592$ & $+0.392932242573$ \\
$R_{\mathrm{ho}}^+$ ($\alpha_1=0.40$) & $+1.076226034357$ & $0.412963868475$ & $+0.621359390645$ & $+1.271224247491$ \\
$R_{\mathrm{ho}}^-$ ($\alpha_2=0.25$) & $+0.787088643179$ & $1.240490193422$ & $-0.454425840016$ & $+0.992253053483$ \\\bottomrule
\end{tabular}
\end{center}
\end{enumerate}
\end{proposition}

\begin{remark}[Physical meaning of the $\ln h$ divergence at $O(\hb^0)$, \Lc/\Ld]
In the full resummed semiclassical amplitude, a one-loop shift of the spectral
parameters $\ell_j^{\mathrm{phys}} = \ell_j + \hb\,\delta\ell_j$ shifts the
resonance locus by $\delta h = O(\hb)$. Expanding
$(\ell_1^{\mathrm{phys}}-\ell_2^{\mathrm{phys}})\ln|\ell_1^{\mathrm{phys}}-\ell_2^{\mathrm{phys}}|/\hb$
at fixed unshifted $h \to 0$ produces $\delta\ell_{12}\ln h$: the $O(\hb^0)$
logarithmic divergence \eqref{eq:delta_renorm_div} is the perturbative signature
that the branch point $h=0$ is shifted by $O(\hb)$ in physical momentum space \Lc/\Ld.
\end{remark}

\subsection{Physical-region singularities and kinematic stability of the light branch}

\begin{theorem}[Kinematic stability of the light branch and heavy-branch threshold embedding, \Lc, \squant]\label{thm:light_stability}
\mbox{}
\begin{enumerate}
\item \textbf{Strict concavity and sub-additivity of $E_{\mathrm{light}}(P)$:}
On the light branch $\alpha \in (0,\pi/3)$ ($|P| \in (0,P_{\max})$ with
$P_{\max} = \sqrt{3}/8$), $E_{\mathrm{light}}(P)$ satisfies $E_{\mathrm{light}}(0)=0$,
$\frac{\dd E_{\mathrm{light}}}{\dd|P|} = \cos\alpha \in (\frac{1}{2},1)$, and
\begin{equation}\label{eq:d2E_dP2}
  \frac{\dd^2 E_{\mathrm{light}}}{\dd P^2} = -\frac{3}{2\sin 3\alpha} < 0
  \qquad \forall\,\alpha \in \left(0,\frac{\pi}{3}\right).
\end{equation}
By strict concavity on $[0,P_{\max}]$ (for same-sign momenta) and strict
monotonicity $|P_1+P_2| < \max(|P_1|,|P_2|)$ (for opposite-sign momenta), any
two nonzero light-branch momenta $P_1,P_2$ with $|P_1+P_2| \le P_{\max}$ obey
\begin{equation}\label{eq:subadditivity}
  \boxed{\;
  E_{\mathrm{light}}(P_1) + E_{\mathrm{light}}(P_2) > E_{\mathrm{light}}(P_1 + P_2).
  \;}
\end{equation}
Consequently, \textbf{$(E,P)$ conservation alone forbids any light soliton from
decaying into two light solitons ($1_{\mathrm{light}}\not\to 2_{\mathrm{light}}$)
and forbids any two light solitons from fusing into a light soliton
($2_{\mathrm{light}}\not\to 1_{\mathrm{light}}$)} \Lc. Since every heavy soliton
has $M = 2k > \sqrt{3}$, $M$-conservation ($k_1+k_2=k_3$) also forbids
$1_{\mathrm{light}}\to\text{heavy}+\text{anything}$. Thus \textbf{light solitons
are kinematically stable against all $1\to 2$ decays} \Lc.

\item \textbf{Every heavy soliton is embedded in the two-light-soliton continuum:}
Conversely, every heavy soliton $S_{13}$ with $\alpha_{13}\in(\pi/3,\pi/2)$ and
direction $\epsilon_{13}\in\{\pm 1\}$ is degenerate in $(E,P,M,\Pi,H_n)$ with the
unique head-on light-light pair $S_{12}$ ($\alpha_{12}=\frac{2\pi}{3}-\alpha_{13}\in(\pi/6,\pi/3)$,
$\epsilon_{12}=+\epsilon_{13}$) and $S_{23}$ ($\alpha_{23}=\alpha_{13}-\frac{\pi}{3}\in(0,\pi/6)$,
$\epsilon_{23}=-\epsilon_{13}$) on $R_{\mathrm{ho}}^+$. The 1+1D two-body
phase-space integral at the heavy soliton's $(E_3,P_3)$ is non-vanishing:
\begin{equation}\label{eq:phase_space_1to2}
  \int \dd P_1\,\dd P_2\,\delta(P_1+P_2-P_3)\,\delta\big(E_1(P_1)+E_2(P_2)-E_3(P_3)\big)
  = \frac{1}{|v_{12} - v_{23}|} = \frac{1}{\sqrt{3}\,k_{13}} \in \left[\frac{1}{\sqrt{3}},\frac{2}{3}\right).
\end{equation}
Across $R_{\mathrm{ho}}^+$ ($\alpha_2 = \frac{\pi}{3}-\alpha_1-h$), the on-shell
energy denominator $\Delta E_{\mathrm{on\text{-}shell}}(h) := E_1(\alpha_1)+E_2(\alpha_2)-E_{\mathrm{heavy}}(P_1+P_2)$
vanishes linearly with non-zero slope:
\begin{equation}\label{eq:onshell_slope}
  \left.\frac{\dd}{\dd h}\Big(E_1(\alpha_1) + E_2(\alpha_2) - E_{\mathrm{heavy}}(P_1+P_2)\Big)\right|_{h=0}
  = -(v_2 - v_3)\,\frac{\dd P_2}{\dd\alpha_2} \ne 0
\end{equation}
(equal to $+0.24107867527655644$ at $\alpha_1=0.62$, verified to $10^{-28}$ by
finite-step differentiation and $10^{-58}$ analytically in \squant) \Lc.
\end{enumerate}
\end{theorem}

\paragraph{Physical-region singularities of the quantum two-body S-matrix (\Lc/\Ld).}
Theorem~\ref{thm:light_stability} and Proposition~\ref{prop:phase_resonance}
distinguish two types of physical-region singularities at the resonant loci:
\begin{enumerate}
\item \textbf{Logarithmic branch point $h^{\ii c h}$ in the WKB/Gamma elastic amplitude \Lc:}
In the semiclassical candidate $S_{\mathrm{cand}} \propto \big(\frac{\Gamma(1+\ii\eta)}{\Gamma(1-\ii\eta)}\big)^\tau\big(\frac{\ee^2 A_{12}}{\eta^2}\big)^{\ii\eta}\ee^{\ii T_{12}/\hb}$
with $\eta = \frac{2}{9\hb}(\ell_1-\ell_2)$, every resonance line has $\eta = 0$,
where $\frac{\Gamma(1+\ii\eta)}{\Gamma(1-\ii\eta)} = 1$ has \textbf{no Gamma pole}
(\texttt{audit\_smx.pdf}). Instead, $A_{12}^{\ii\eta} \sim h^{\ii\sigma_{\mathrm{res}}c_\ell h/(9\hb/2)}$
produces a continuous \textbf{logarithmic branch point} at real rapidity $h=0$,
across which $A_{12}$ changes sign \Lc.
\item \textbf{Real-rapidity $s$-channel pole / resonance threshold from $2\leftrightarrow 1$ coupling \Ld:}
Because $\Delta E_{\mathrm{on\text{-}shell}}(h) = E_1+E_2-E_{\mathrm{heavy}}(P_1+P_2)$
vanishes linearly on $R_{\mathrm{ho}}^+$ \eqref{eq:onshell_slope}, any non-zero
quantum $2\leftrightarrow 1$ transition vertex $g_{12\to 3}$ generates an
on-shell $s$-channel propagator pole $\big(E_1+E_2-E_3(P_1+P_2) - \Sigma_3(E,P) + \ii\epsilon\big)^{-1}$
directly on the real head-on scattering domain $R_{\mathrm{ho}}^+$ (and crossed
$t/u$-channel singularities on $R_{\mathrm{ho}}^-$ and $R_{\mathrm{ov}}$) \Ld.
By Dyson resummation of the two-light-soliton bubble \eqref{eq:phase_space_1to2},
the embedded heavy soliton $S_{13}$ acquires a decay width
$\Gamma_{13\to 12+23} = \frac{|g_{12\to 3}|^2}{2\sqrt{3}\,k_{13}} \ge 0$, turning
the real pole into a Breit--Wigner resonance just below the real axis \Ld.
\end{enumerate}

\subsection{Two-soliton unitarity vs.\ coupled \texorpdfstring{$2\leftrightarrow 1$}{2<->1} channels}

\begin{proposition}[WKB unitarity violation in the singular domain $A_{12}<0$, \Lc, \squant]\label{prop:unitarity_deficit}
When the semiclassical phase \eqref{eq:delta_semi} (or the continued Gamma
candidate) is continued across any of $R_{\mathrm{ov}},R_{\mathrm{ho}}^+,
R_{\mathrm{ho}}^-$ into the singular region $A_{12} < 0$, $\ln A_{12} = \ln|A_{12}| \pm \ii\pi$
acquires an imaginary part $\pm\ii\pi$. On the physical sheet where the elastic
amplitude remains bounded ($\operatorname{Im}\delta > 0$),
\begin{equation}\label{eq:wkb_deficit}
  \boxed{\;
  \big|S_{2\to 2}^{\mathrm{WKB}}(1,2)\big|^2
  = \left|\ee^{\ii\delta(1,2)/\hb}\right|^2
  = \exp\!\left(-\frac{4\pi}{9\hb}\,\big|\ell_1 - \ell_2\big|\right) < 1
  \qquad (A_{12} < 0).
  \;}
\end{equation}
For example, at the singular head-on point $(\alpha_1,\alpha_2)=(0.50,0.65)$
($\epsilon_1=+1,\epsilon_2=-1$, where $A_{\mathrm{ho}} = -0.354784 < 0$ and
$|\ell_1-\ell_2| = 0.295758$) with $\hb=0.1$, $|S_{2\to 2}^{\mathrm{WKB}}|^2 = 0.0160837 < 1$,
leaving a two-body unitarity deficit $1 - |S_{2\to 2}^{\mathrm{WKB}}|^2 = 0.9839163 > 0$ \Lc.
\end{proposition}

\paragraph{Consequence for unitarity (\Lc/\Ld).}
\begin{itemize}
\item \textbf{At the resonance curve $R_{\mathrm{ho}}^+$ ($E_1+E_2=E_3(P_1+P_2)$):}
Because the two-light-soliton state $|P_1,P_2\rangle$ and the heavy one-soliton
state $|P_3=P_1+P_2\rangle$ share identical conserved charges, unitarity on the
energy-momentum shell acts on the coupled two-dimensional space
$\mathcal{H}_{E,P} \simeq \mathbb{C}|S_{12},S_{23}\rangle \oplus \mathbb{C}|S_{13}\rangle$ via a
$2\times 2$ unitary block
\begin{equation}\label{eq:coupled_S}
  \mathbf{S}(E,P) = \begin{pmatrix} S_{2\to 2} & S_{1\to 2} \\ S_{2\to 1} & S_{1\to 1} \end{pmatrix}, \qquad
  |S_{2\to 2}|^2 + |S_{2\to 1}|^2 = 1.
\end{equation}
Two-soliton unitarity alone ($|S_{2\to 2}|=1,\,S_{2\to 1}=0$) is only consistent
if the quantum $2\leftrightarrow 1$ transition amplitude vanishes identically on
$R_{\mathrm{ho}}^+$ despite the existence of the classical Y-junction trajectory
connecting the two states \Ld.
\item \textbf{In the singular region $A_{12}<0$:}
Pure two-soliton elastic unitarity fails at the semiclassical level by
\eqref{eq:wkb_deficit} \Lc. Because the classical 4-exponential solution blows up
($\tau=0$) whenever $A_{12}<0$, two incoming solitons with $A_{12}<0$ cannot pass
through one another as a regular classical two-soliton state; in the quantum
theory \Ld, the absorption deficit $1 - \exp(-\frac{4\pi}{9\hb}|\ell_1-\ell_2|)$
measures the suppression of direct elastic transmission in favour of
inelastic/multi-soliton/radiation or resonant rearrangement channels.
\end{itemize}

\subsection{Three-soliton factorisation and bulk fusion bootstrap near resonance}

In 1+1D diagonal scattering, the algebraic Yang--Baxter equation
$S_{12}S_{13}S_{23} = S_{23}S_{13}S_{12}$ is trivial, but physical
three-body factorisation requires that a three-soliton collision decompose into
independent pairwise events regardless of chronological order, and---when two
solitons fuse ($S_{12}+S_{23}\to S_{13}$)—that scattering a third soliton $3$
off the pair before fusion equal scattering off the fused soliton $S_{13}$ after
fusion (the bulk fusion bootstrap).

\begin{proposition}[Classical and bare 1-loop fusion factorisation vs.\ $O(\hb^0)$ obstructions, \Lc, \squant]\label{prop:three_body}
Let $S_{12}=(q_1,q_2)$, $S_{23}=(q_2,q_3)$, and $S_{13}=(q_1,q_3)$ be the three
solitons of a KP triple ($q_1>q_2>q_3$, sharing $\ell_{\mathrm{res}}$), and let
$3 = (\alpha_3,\epsilon_3)$ be any test soliton in the regular region
($A(3,S_{ij}) > 0$).
\begin{enumerate}
\item \textbf{Exact classical cross-ratio and Darboux-phase fusion:}
The shared intermediate root $q_2$ cancels identically in the product of KP
cross-ratios,
\begin{align}\label{eq:A_fusion}
  A(3,S_{12})\,A(3,S_{23})
  &= \frac{(p_{3a}-q_1)(p_{3b}-q_2)}{(p_{3a}-q_2)(p_{3b}-q_1)}\cdot
     \frac{(p_{3a}-q_2)(p_{3b}-q_3)}{(p_{3a}-q_3)(p_{3b}-q_2)} \nonumber\\
  &= \frac{(p_{3a}-q_1)(p_{3b}-q_3)}{(p_{3a}-q_3)(p_{3b}-q_1)}
   = A(3,S_{13}),
\end{align}
and because $M(S_{12})+M(S_{23})=M(S_{13})$, $\Pi(S_{12})+\Pi(S_{23})=\Pi(S_{13})$,
and $\ell(S_{12})=\ell(S_{23})=\ell(S_{13})=\ell_{\mathrm{res}}$, the full
Darboux-frame semiclassical phase \eqref{eq:delta_semi} is \textbf{strictly
additive across bulk fusion with zero topological remainder}:
\begin{equation}\label{eq:delta_bulk_fusion}
  \boxed{\;
  \delta_{\mathrm{cl}}(3,S_{12}) + \delta_{\mathrm{cl}}(3,S_{23})
  = \delta_{\mathrm{cl}}(3,S_{13}).
  \;}
\end{equation}
Thus the pure WKB exponential $S = \ee^{\ii\delta_{\mathrm{cl}}/\hb}$ obeys the
bulk fusion bootstrap $S(3,S_{12})\,S(3,S_{23}) = S(3,S_{13})$ identically \Lc.

\item \textbf{Pointwise phonon phase telescoping and bare one-loop integral fusion:}
For every phonon wavenumber $q \in \mathbb{R}\setminus\{0\}$, the one-soliton
phonon phase derivatives telescope across the KP triple,
\begin{equation}\label{eq:phonon_dt_telescope}
  \theta_{\mathrm{sol}}'(q;S_{12}) + \theta_{\mathrm{sol}}'(q;S_{23})
  = \theta_{\mathrm{sol}}'(q;S_{13}),
\end{equation}
and therefore the bare fixed-rapidity one-loop continuum phase integrals satisfy
$I_{SS}(3,S_{12}) + I_{SS}(3,S_{23}) = I_{SS}(3,S_{13})$ identically \Lc.

\item \textbf{Two $O(\hb^0)$ breakdowns of bulk fusion factorisation:}
\begin{itemize}
\item \emph{Maslov/Stirling mismatch in the continued Gamma candidate:}
Because $S_{12}$ ($\epsilon_{12}=+1$) and $S_{23}$ ($\epsilon_{23}=-1$) move in
opposite directions, their relative orientation signs $\tau(3,S_{12})+\tau(3,S_{23}) = (+1)+(-1) = 0$
cancel, whereas $\tau(3,S_{13}) = \pm 1 \ne 0$. With common
$\eta = \frac{2}{9\hb}(\ell_3-\ell_{\mathrm{res}})$, the continued Gamma candidate
violates bulk fusion at $O(\hb^0)$ by:
\begin{align}\label{eq:gamma_fusion_fail}
  \frac{S_{\mathrm{cand}}(3,S_{12})\,S_{\mathrm{cand}}(3,S_{23})}{S_{\mathrm{cand}}(3,S_{13})}
  &= \left[\frac{\Gamma(1+\ii\eta)}{\Gamma(1-\ii\eta)}\!\left(\frac{\ee^2}{\eta^2}\right)^{-\ii\eta}\right]^{-\tau(3,S_{13})} \nonumber\\
  &= \ee^{-\ii\frac{\pi}{2}\tau(3,S_{13})\sgn\eta}\big(1 + O(\hb)\big) \ne 1
\end{align}
(at $\phi=0.23,\alpha_3=0.31,\hb=0.1$, the phase discrepancy is
$+1.5158110455271881045$ vs.\ $+\frac{\pi}{2} = +1.5707963267948966192$ as
$\hb\to 0^+$) \Lc.
\item \emph{Fixed-physical-momentum target-soliton renormalisation mismatch:}
When $S(3,S_{ij})$ is evaluated at fixed physical momenta
$(P_3^{\mathrm{phys}},P_{S_{ij}}^{\mathrm{phys}})$, the test-soliton Taylor shifts
$\Delta X_3(3,S_{ij})\Delta P_3$ cancel by \eqref{eq:A_fusion}, but the
target-soliton Taylor shifts fail to telescope because $\Delta P(S_{ij})/k_{S_{ij}}$
varies across the three chords:
\begin{equation}\label{eq:target_P_mismatch}
  \sum_{S\in\{S_{12},S_{23}\}} \frac{\ln A(3,S)}{k_S}\Delta P(S)
  - \frac{\ln A(3,S_{13})}{k_{S_{13}}}\Delta P(S_{13}) \ne 0
\end{equation}
(equal to $+0.8755292547543045976$ at $\phi=0.23,\alpha_3=0.31$) \Lc.
\end{itemize}
\end{enumerate}
\end{proposition}

\begin{remark}[Wave-packet / Landau breakdown of 3-body chronological independence near resonance, \Ld]
Classically, \eqref{eq:A_fusion}--\eqref{eq:delta_bulk_fusion} hold because the
fused intermediate soliton $S_{13}$ carries the exact sum of the KP chords of
$S_{12}$ and $S_{23}$. Quantum mechanically \Ld, near the resonance locus
($|h|\ll 1$), the divergent time delay $\Delta t_{12} \sim \frac{1}{|v_1-v_2|k_1}\ln(1/|h|)$
means that in a 3-body wave-packet collision with impact parameter separation
$\Delta X \sim v\,\Delta t_{12}$, the intermediate state is an on-shell coherent
superposition of $|S_{12},S_{23}\rangle$ and $|S_{13}\rangle$. Because the $O(\hb^0)$
scattering phases off $|S_{12},S_{23}\rangle$ vs.\ $|S_{13}\rangle$ differ by the
nonzero target-momentum renormalisation \eqref{eq:target_P_mismatch} (unless
cancelled by a two-body quantum frame counterterm), chronological independence of
$3\to 3$ scattering fails at $O(\hb^0)$ in the resonant transition zone \Ld.
\end{remark}

% ==============================================================================
\section{Question 4: Conclusion---when is the quantum S-matrix elastic and factorised?}
\label{sec:conclusion}
% ==============================================================================

Combining our exact structural results \Lc\ with their quantum implications \Ld,
we answer Question~4 by partitioning the two-soliton kinematic space into four
physically distinct regimes:

\begin{enumerate}
\item \textbf{Regime I: Light-light overtaking collisions ($\alpha_1,\alpha_2\in(0,\pi/3),\,\epsilon_1=\epsilon_2$).}\\
\emph{Status:} \textbf{Elastic and factorised} (\Lc\ for kinematic protection, \Ld\ for full quantum theory).
\begin{itemize}
\item Every soliton has $|P_j| < P_{\max} = \sqrt{3}/8$ on the KdV-like light
branch, $\alpha_1+\alpha_2 < 2\pi/3$ everywhere, so $0 < A_{\mathrm{ov}} < 1$ is
strictly regular and bounded away from $R_{\mathrm{ov}}$ except at the double-cusp
corner $(\pi/3,\pi/3)$ \Lc.
\item By Theorem~\ref{thm:light_stability}(1), $1_{\mathrm{light}}\to 2$ decay and
$2_{\mathrm{light}}\to 1_{\mathrm{light}}$ fusion are strictly forbidden by
$(E,P)$ conservation, and $2_{\mathrm{light}}\to 1_{\mathrm{heavy}}$ overtaking
fusion is forbidden by $(M,P,E)$ conservation ($\alpha_1+\alpha_2 < 2\pi/3$) \Lc.
\item No resonant or number-changing classical channel is reachable; a unitary,
diagonal, factorised quantum two-soliton S-matrix is fully consistent here \Ld.
\end{itemize}

\item \textbf{Regime II: Sub-threshold light-light head-on collisions ($\alpha_1+\alpha_2 < \pi/3,\,\epsilon_1=-\epsilon_2$).}\\
\emph{Status:} \textbf{Elastic below the heavy-soliton threshold $R_{\mathrm{ho}}^+$} (\Lc/\Ld).
\begin{itemize}
\item In this triangle ($\frac{1}{2}$ of the light-light square, $\frac{2}{9}$
of $Q$), $A_{\mathrm{ho}} > 1$ is strictly regular and attractive \Lc.
\item The total energy lies strictly below the heavy-soliton threshold:
$E_1(\alpha_1) + E_2(\alpha_2) < E_{\mathrm{heavy}}(P_1+P_2)$ (since
$\Delta E_{\mathrm{on\text{-}shell}}(h) > 0$ for $h=\pi/3-(\alpha_1+\alpha_2)<0$
and $<0$ for $h>0$ when $\alpha_1>\pi/6$, or vanishes only on $R_{\mathrm{ho}}^+$),
and $(M,P,E)$ conservation forbids $2\to 1$ fusion anywhere off $R_{\mathrm{ho}}^+$ \Lc.
\item As $\alpha_1+\alpha_2 \to (\pi/3)^-$, the collision approaches the
on-shell heavy-soliton resonance $R_{\mathrm{ho}}^+$, where the classical delay
diverges as $\ln(1/h)$ and the $s$-channel heavy-soliton pole/branch point sits
on the boundary of the elastic domain \Lc/\Ld.
\end{itemize}

\item \textbf{Regime III: Heavy-soliton channels ($\alpha_j \in (\pi/3,\pi/2)$) and the resonant loci $R_{\mathrm{ho}}^+,R_{\mathrm{ho}}^-,R_{\mathrm{ov}}$.}\\
\emph{Status:} \textbf{Inelastic $1\leftrightarrow 2$ and resonant rearrangement channels open} (\Lc\ for kinematics/conservation, \Ld\ for quantum amplitudes).
\begin{itemize}
\item Every heavy soliton $S_{13}$ ($\alpha_{13}\in(\pi/3,\pi/2)$) is an embedded
state in the two-light-soliton head-on continuum with exact conservation of
$M,\Pi,P,E$ and all $H_n$ into $S_{12}+S_{23}$ and non-zero 1+1D phase space
$J_{1\to 2}=(\sqrt{3}\,k_{13})^{-1}$ \Lc. Unless a hidden quantum selection rule
forces the $1\leftrightarrow 2$ matrix element to vanish identically, heavy
solitons are unstable quantum resonances ($S_{13}\to S_{12}+S_{23}$) rather than
stable asymptotic particles \Ld.
\item Semiclassically, near $R_{\mathrm{ho}}^+$ ($h = \pi/3 - (\alpha_1+\alpha_2) \to 0$),
the $2\to 2$ amplitude couples to the $2\to 1$ fusion channel via the on-shell
energy denominator $\Delta E_{\mathrm{on\text{-}shell}}(h) \propto h$, while
across the resonance curve ($h < 0$, entering $A_{12} < 0$) the direct WKB
elastic transmission amplitude is suppressed by the barrier-penetration factor \Lc/\Ld:
\begin{equation}\label{eq:inelastic_amp}
  \big|S_{2\to 2}^{\mathrm{WKB}}\big| = \exp\!\left(-\frac{2\pi}{9\hb}|\ell_1-\ell_2|\right), \qquad
  P_{\mathrm{inelastic}}^{\mathrm{WKB}} = 1 - \big|S_{2\to 2}^{\mathrm{WKB}}\big|^2
  = 1 - \exp\!\left(-\frac{4\pi}{9\hb}|\ell_1-\ell_2|\right).
\end{equation}
\end{itemize}

\item \textbf{Regime IV: Open singular two-soliton domains ($A_{12} < 0$) and complex Hirota breathers ($A_{\mathrm{br}} < 1/4$).}\\
\emph{Status:} \textbf{Non-elastic / outside the regular soliton sector} (\Lc/\Ld).
\begin{itemize}
\item For $\alpha_1+\alpha_2 > 2\pi/3$ (overtaking) and $\alpha_1+\alpha_2 > \pi/3,\,|\alpha_1-\alpha_2| < \pi/3$
(head-on), $A_{12} < 0$: the classical 4-exponential solution develops $\tau=0$
poles, and the continued WKB elastic amplitude violates $|S_{2\to 2}|=1$ by
\eqref{eq:wkb_deficit} \Lc. Even in the light-light head-on square $(0,\pi/3)^2$,
the upper triangle $\alpha_1+\alpha_2 \in (\pi/3,2\pi/3)$ lies in $A_{\mathrm{ho}} < 0$,
so \textbf{even restricting external states to the light branch $\alpha<\pi/3$
does not make head-on scattering globally elastic for $\alpha_1+\alpha_2 > \pi/3$} \Lc/\Ld!
\item Every two-phase zero-background Hirota breather of the good equation has
$A_{\mathrm{br}} < 1/4 < 1$ and is classically singular (\texttt{phase1/\allowbreak review.pdf},
\texttt{audit\_smx.pdf}), so there is no regular classical breather family to
require Bohr--Sommerfeld bound-state poles on the physical sheet \Lc/\Ld.
\end{itemize}
\end{enumerate}

% ==============================================================================
\section{Summary of claims, labels, and verification scripts}
\label{sec:summary_table}
% ==============================================================================

All three verification scripts in \texttt{phase2/\allowbreak scripts/} run in
under 25~seconds total and terminate with \texttt{ALL CHECKS PASSED}.

\small
\begin{longtable}{@{}L{4.8cm}L{1.2cm}L{4.3cm}L{4.2cm}@{}}\toprule
Claim / Formula & Label & Script & Precision / Check \\\midrule\endhead
KP chord dictionary \eqref{eq:kp_roots}--\eqref{eq:A12_crossratio}, cross-ratio $=$ Hirota $=$ trig forms & \La/\Lc & \skin & Exact symbolic over $\mathbb{Q}(\sqrt{3})[t_1,t_2]$ \\
Exact resonance lines $R_{\mathrm{ov}},R_{\mathrm{ho}}^+,R_{\mathrm{ho}}^-$; refutation of $|\alpha_1-\alpha_2|=\pi/3$ for overtaking & \Lc & \skin & 60-digit mpmath ($<10^{-48}$) \\
KP-triangle unification \eqref{eq:triple_chords}--\eqref{eq:three_projections}: 2 light chords $+$ 1 heavy chord & \Lc & \skin & Symbolic $+$ 60-digit mpmath \\
Exact 3-exponential Y-junction PDE solution and $1\leftrightarrow 2$ asymptotic limits & \La/\Lc & \skin & 60-digit PDE jet ($<10^{-50}$) and $t=\pm 400$ ($<10^{-48}$) \\
Kinematic area fractions ($\frac{7}{9}$ reg.\ ov., $\frac{1}{3}$ reg.\ ho., $100\%$ light-light ov.) & \Lc & \skin & Exact rational geometry \\
Symbolic integration of $H_1,\dots,H_8$ via $T=\tanh(kx/2)$ \eqref{eq:H1_to_H8} & \Lc & \scons & Exact symbolic polynomial integration \\
Univariate chord potentials $F_Q(p)$ \eqref{eq:F_potentials} for $M,\Pi,P,E$ and $H_1,\dots,H_8$ & \Lc & \scons & Exact polynomial division mod $p_a^2+p_a p_b+p_b^2-\frac{1}{4}$ \\
Spin-3 gap: pure $k^3 \notin \operatorname{span}\{\Delta_1,\Delta_5\}$, while $E\in\operatorname{span}\{\Delta_1,\Delta_5\}$ & \Lc & \scons & Exact symbolic linear system \\
Classical telescoping $Q(S_{12})+Q(S_{23})=Q(S_{13})$ and Y-junction quadrature across $t\in\{-2,0,2\}$ & \Lc & \scons & Exact symbolic $+$ 60-digit quadrature ($<10^{-50}$) \\
One-loop mode-sum telescoping vs.\ fixed-$P$ mismatch $\Delta(F|_P)=0.3887924407\ldots$ \eqref{eq:delta_F_at_P} & \Lc & \scons & 60-digit mpmath ($<10^{-48}$) \\
Transverse $h\ln h$ branch point \eqref{eq:delta_hlnh} and fixed-$P$ $\ln A_{12}$ divergence \eqref{eq:delta_renorm_div} & \Lc & \squant & 60-digit asymptotic extraction ($<10^{-18}$) \\
Strict concavity $\frac{\dd^2 E_{\mathrm{light}}}{\dd P^2}=-\frac{3}{2\sin 3\alpha}<0$ and $1_{\mathrm{light}}\not\leftrightarrow 2_{\mathrm{light}}$ \eqref{eq:subadditivity} & \Lc & \squant & Exact symbolic $+$ 60-digit grid check \\
Heavy-branch $1\to 2$ Jacobian \eqref{eq:phase_space_1to2} and linear on-shell energy zero \eqref{eq:onshell_slope} & \Lc & \squant & 60-digit analytic $+$ finite-step ($<10^{-28}$) \\
WKB unitarity deficit $|S_{2\to 2}^{\mathrm{WKB}}|^2 = \ee^{-\frac{4\pi}{9\hb}|\ell_1-\ell_2|} < 1$ for $A_{12}<0$ \eqref{eq:wkb_deficit} & \Lc & \squant & 60-digit evaluation \\
Classical bulk fusion $A_{3,12}A_{3,23}=A_{3,13}$ and $\delta_{\mathrm{cl}}(3,S_{12})+\delta_{\mathrm{cl}}(3,S_{23})=\delta_{\mathrm{cl}}(3,S_{13})$ & \Lc & \squant & 60-digit mpmath ($<10^{-48}$) \\
Pointwise phonon telescoping \eqref{eq:phonon_dt_telescope}; Gamma candidate $\ee^{\mp\ii\pi/2}$ and target-$\Delta P$ obstructions & \Lc & \squant & 60-digit mpmath ($<10^{-48}$) \\
Physical elasticity conditions, heavy-soliton decay width, and inelastic regime partition & \Ld & --- & Conditional on quantum vertex and frame map \\\bottomrule
\end{longtable}
\normalsize

% ==============================================================================
\section{Failures recorded during development}
\label{sec:failures}
% ==============================================================================

In accordance with \texttt{AGENTS.md}, we record every failed check or error
encountered during the development of this package and how it was resolved
without weakening any check:

\begin{enumerate}
\item \textbf{SymPy trigonometric cube simplification and rational GCD hang in \skin:}
In the initial version of \skin, \texttt{sp.simplify(sp.expand\_trig(4*p\_root**3 - p\_root - ell1/(3*sp.sqrt(3))))}
returned a non-zero unexpanded expression containing \texttt{sin(alpha\_1 + pi/6)**3},
and \texttt{sp.cancel} over $\mathbb{Q}(\sqrt{3})[t_1,t_2]$ hung on rational
cross-ratio differences. Clearing polynomial denominators using the homogeneous
Weierstrass parametrization $S_j = 2t_j,\,C_j = 1-t_j^2,\,D_j = 1+t_j^2$
reduced all identities to pure polynomial expansions in $\mathbb{Q}(\sqrt{3})[t_1,t_2]$
that vanish identically in $<0.1\,\mathrm{s}$.

\item \textbf{Trailing-digit fabrication on the false overtaking resonance value in \skin, and insufficient asymptotic time $T=60$:}
In \skin, an initial assertion compared $A_{\mathrm{ov}}(0.2,0.2+\pi/3)$ to a
20-digit string obtained by appending arbitrary digits to the 12-digit number
$0.355465479849$ quoted in \texttt{audit\_smx.tex}; the assertion failed because
the true 25-digit value is $0.3554654798490616496187723$. It was replaced by an
exact $10^{-48}$ check against the closed form $\frac{1-3\cot^2(\alpha_1+\pi/6)}{-8}$
together with a $10^{-12}$ check against the rounded audit figure. In the same
script, checking the $t\to\pm\infty$ peak height of the Y-junction at $T=60$
failed the $10^{-25}$ tolerance because the transverse exponent
$\sqrt{3}k_{12}k_{23}T \approx 17.3$ gives $\ee^{-17.3}\approx 3\times 10^{-8}$;
evaluating at $T=400$ ($\ee^{-115} < 10^{-49}$) passed the stricter $10^{-48}$
tolerance.

\item \textbf{Initial mental guess for the rational coefficients of $H_7$ and $H_8$ in \scons:}
Before running the symbolic integrator in \scons, the docstring contained a
tentative hand-estimated decomposition for $H_7$ and $H_8$. Exact symbolic
integration of $\Res L^{7/3}$ and $\Res L^{8/3}$ via $T=\tanh(kx/2)$ yielded the
unique exact coefficients $H_7 = \Delta_7 - \frac{7}{12}\Delta_5 + \frac{7}{324}\Delta_1$
and $H_8 = \Delta_8 - \frac{7}{36}\Delta_4 + \frac{11}{324}\Delta_2$, which
pass exact polynomial division modulo $p_a^2+p_a p_b+p_b^2-\frac{1}{4}$.

\item \textbf{Unbracketed Newton overshoot across the cusp $\alpha=\pi/3$, sign of $P_1+P_2$ on $R_{\mathrm{ho}}^+$, and an extra $\sin\alpha$ factor in $\dd P/\dd\alpha$ in \squant:}
In the first run of \squant, inverting $P(\alpha) = P_{\mathrm{target}}$ on the
light branch by unbracketed Newton iteration from $\alpha_0=0.5$ overshot the
cusp $\alpha=\pi/3$ (where $P'(\pi/3)=0$) onto the heavy branch when
$P_1+P_2$ was near $P_{\max}$, causing the sub-additivity assertion
$E_{\mathrm{light}}(P_1)+E_{\mathrm{light}}(P_2) > E_{\mathrm{light}}(P_1+P_2)$
to fail. Replacing unbracketed Newton by bracketed bisection on $[0,\pi/3]$
(light) and $[\pi/3,\pi/2]$ (heavy) followed by Newton polishing resolved the
branch jump. In the same test, choosing $\alpha_1=0.42 < \pi/6$ ($\epsilon_1=+1$)
and $\alpha_2=\pi/3-0.42 > \pi/6$ ($\epsilon_2=-1$) gave $P_1+P_2 < 0$, so the
fused heavy soliton had $\epsilon_3 = -1$ rather than $+1$; and the Newton
polishing step initially wrote $\frac{\dd P}{\dd\alpha} = \frac{2}{3}\sin^2\alpha\sin 3\alpha$
instead of the correct $\frac{2}{3}\sin\alpha\sin 3\alpha$, which degraded the
finite-step derivative check of $\Delta E_{\mathrm{on\text{-}shell}}'(0)$.
Correcting $\frac{\dd P}{\dd\alpha}$ and testing $\alpha_1=0.62 \in (\pi/6,\pi/3)$
($P_1+P_2 > 0,\,\epsilon_3=+1$) passed at $10^{-28}$ (finite-step) and $10^{-48}$
(analytic).
\end{enumerate}

% ==============================================================================
\paragraph{Reference verification and reading scope.}
All references below are inherited from the verified repository bibliography
(\texttt{phase2/\allowbreak report.tex} and \texttt{phase2/\allowbreak errata.tex});
no new external reference is introduced. Reading scope follows the project
ledger: relevant full-text sections read in the project for \cite{BC,BD,CZ,DMHC,FK,RS};
abstract/metadata verified in \texttt{errata.tex} for \cite{Dickey,HI,LMK,Zak}.

\begin{thebibliography}{99}\footnotesize\setlength{\itemsep}{1pt}
\bibitem{BC} M.~M.~Bogdan, O.~V.~Charkina, \emph{Quasiclassical quantization of the Boussinesq breather emerging from the kink localized mode}, \emph{Low Temp. Phys.} 51 (2025) 763--773, arXiv:2502.14695, doi:\url{10.1063/10.0036799}. [relevant full-text passages read; bad equation]
\bibitem{BD} T.~Bonnemain, B.~Doyon, \emph{Soliton gas of the integrable Boussinesq equation and its generalised hydrodynamics}, \emph{SciPost Phys.} 18 (2025) 075, arXiv:2402.08669. [full text read]
\bibitem{CZ} E.~Corrigan, C.~Zambon, \emph{Adding integrable defects to the Boussinesq equation}, \emph{J. Phys. A: Math. Theor.} 56 (2023) 385701, arXiv:2306.03772, doi:\url{10.1088/1751-8121/aceec9}. [full text read]
\bibitem{Dickey} L.~A.~Dickey, \emph{Soliton Equations and Hamiltonian Systems}, 2nd ed., World Scientific (2003), doi:\url{10.1142/5108}. [metadata verified; not read in full]
\bibitem{DMHC} A.~Dimakis, F.~M\"uller-Hoissen, X.-M.~Chen, \emph{ Soliton solutions of the matrix Boussinesq equation and beyond Yang--Baxter maps}, \emph{Phys. Scr.} 94 (2019) 035206, arXiv:1805.09711, doi:\url{10.1088/1402-4896/aaf6da}. [full text read]
\bibitem{FK} L.~D.~Faddeev, V.~E.~Korepin, \emph{Quantum theory of solitons}, \emph{Phys. Rep.} 42 (1978) 1--87, doi:\url{10.1016/0370-1573(78)90058-3}. [full text read]
\bibitem{HI} R.~Hirota, M.~Ito, \emph{Resonance of solitons in one dimension}, \emph{J. Phys. Soc. Jpn.} 52 (1983) 744--748, doi:\url{10.1143/JPSJ.52.744}. [abstract/metadata only]
\bibitem{LMK} F.~Lambert, M.~Musette, E.~Kesteloot, \emph{Soliton resonances for the good Boussinesq equation}, \emph{Inverse Problems} 3 (1987) 275--288, doi:\url{10.1088/0266-5611/3/2/010}. [abstract/metadata only]
\bibitem{RS} A.~G.~Rasin, J.~Schiff, \emph{B\"acklund transformations for the Boussinesq equation and merging solitons}, \emph{J. Phys. A: Math. Theor.} 50 (2017) 325202, arXiv:1701.00243, doi:\url{10.1088/1751-8121/aa7942}. [full text read]
\bibitem{Zak} V.~E.~Zakharov, \emph{On sto\-chastic\-i\-za\-tion of one-dimensional chains of nonlinear oscillators}, \emph{Sov. Phys. JETP} 38 (1974) 108--110. [abstract/citation record only]
\end{thebibliography}

\end{document}
